% =============================================================================
% theorems.tex - Probability Notes theorem system and section skeleton (v1.1)
% =============================================================================
% Input after shared/components/math.tex. Loads amsthm only.
\RequirePackage{amsthm}

\newtheorem{theorem}{Theorem}
\newtheorem{proposition}[theorem]{Proposition}
\newtheorem{lemma}[theorem]{Lemma}
\newtheorem{corollary}[theorem]{Corollary}
\theoremstyle{definition}
\newtheorem{definition}[theorem]{Definition}
\newtheorem{assumption}{Assumption}
\renewcommand{\theassumption}{A\arabic{assumption}}
\newtheorem{example}[theorem]{Example}
\theoremstyle{remark}
\newtheorem{remark}[theorem]{Remark}
\renewcommand{\qedsymbol}{\textcolor{Primary}{$\square$}}

% cleveref names. The class loads cleveref before amsthm, so shared-counter
% environments must be labelled with their type: \label[lemma]{lem:x},
% \label[corollary]{cor:x}, \label[remark]{rem:x}. Theorems use plain \label.
\crefname{theorem}{Theorem}{Theorems}
\crefname{lemma}{Lemma}{Lemmas}
\crefname{corollary}{Corollary}{Corollaries}
\crefname{assumption}{Assumption}{Assumptions}
\crefname{remark}{Remark}{Remarks}

% Proof steps and flagged interchanges (\justify is taken by ragged2e)
\newcommand{\step}[1]{\par\smallskip\noindent\textbf{\color{PrimaryDark}Step: #1.}\ }
\newcommand{\interchange}[1]{\par\smallskip\noindent\textit{Justification (#1).}\ }

% Mandatory section skeleton: use all eight, in this order.
\newcommand{\SecSetup}{\section{Setup and notation}}
\newcommand{\SecResult}{\section{Result}}
\newcommand{\SecProof}{\section{Proof}}
\newcommand{\SecInsight}{\section{Insights}}
\newcommand{\SecUses}{\section{Uses and consequences}}
\newcommand{\SecExample}{\section{Worked example}}
\newcommand{\SecCheck}{\section{Numerical check}}
\newcommand{\SecNotes}{\section{Notes}}

% Headline result: exactly one per issue, never split across pages.
\newtcolorbox{mainresult}[1][Main result]{saim/green,title={#1},breakable=false}
